\documentclass[11pt]{article}
\usepackage[letterpaper,margin=1.05in]{geometry}
\usepackage[T1]{fontenc}
\usepackage{lmodern,amsmath,amssymb,amsthm,microtype,booktabs}
\usepackage[colorlinks=true,linkcolor=blue,citecolor=blue,urlcolor=blue]{hyperref}
\newtheorem{theorem}{Theorem}
\newtheorem{lemma}[theorem]{Lemma}
\newtheorem{corollary}[theorem]{Corollary}
\newtheorem{proposition}[theorem]{Proposition}
\theoremstyle{definition}
\newtheorem{remark}[theorem]{Remark}
\newcommand{\Om}{\widetilde\Omega}
\newcommand{\F}{\widetilde{\mathcal F}}
\newcommand{\D}{\widetilde{\mathcal D}}
\newcommand{\Rp}{\mathcal R_{\mathrm{post}}}
\newcommand{\Pinfty}{\mathcal P_\infty}
\newcommand{\cyl}{\operatorname{cyl}}
\newcommand{\TV}{\operatorname{TV}}
\title{Complex percolation on the doubly-principal\protect\\ cylinder sigma-algebra}
\author{Working manuscript}
\date{21 October 2026}
\begin{document}
\maketitle
\begin{abstract}
For complex transitive percolation, a countably additive complex measure
matching the prescribed naturally labelled doubly-principal cylinder
amplitudes exists exactly when the percolation parameter belongs to $[0,1]$.
The obstruction already appears on the sigma-algebra generated by these
cylinders, without using all finite cylinders. Two universal caps reduce the
absolute level sum to a weighted sum over arbitrary finite orders.
Normalizing the modulus weights gives a real random graph order and its
standard endpoint-connectivity product, yielding an exponentially divergent
lower bound outside the real segment. This completes a calculation suggested
by Dowker--Zalel, rather than introducing a general extension criterion. A separate covariant result requires an all-principal common-factor
rule and normalization on infinitely many posts. A Dirac counterexample
shows why that rule cannot be weakened to doubly-principal pasts alone.
\end{abstract}

\section{The labelled classification}
Let $[n]=\{0,\ldots,n-1\}$. A naturally labelled finite order $C$ on $[n]$
has $i\prec_C j$ only if $i<j$. Write $L(C)$ for its number of covers,
$R(C)$ for its number of strict comparable pairs, and
$h(C)=\binom n2-R(C)$ for its number of incomparable pairs. For $p\in\mathbb C$
and $q=1-p$, prescribe
\begin{equation}\label{eq:amp}
 A_p(C)=p^{L(C)}q^{h(C)}.
\end{equation}
Powers with exponent zero equal one, also at $p=0$ or $q=0$.

The space $\Om$ consists of orders on $\mathbb N$ with the same natural-label
property. Its usual sigma-algebra $\F$ is generated by finite-prefix cylinders
$\cyl(C)=\{\omega:\omega|_{[n]}=C\}$. An order is \emph{principal} if it has
a unique maximal element. For $n\ge2$, it is \emph{doubly principal} if both
$C$ and $C|_{[n-1]}$ are principal. Set
\begin{equation}\label{eq:algebra}
 \D=\sigma\{\cyl(C): C\text{ is doubly principal}\}\subseteq\F.
\end{equation}
We write $\D$ to avoid using a numerical parameter as the subscript for an
algebra that is independent of $p$. This is the labelled algebra denoted
$\widetilde{\mathcal R}_p$ in Dowker--Zalel~\cite[Sec.~V]{DZ}.

\begin{theorem}[Labelled doubly-principal classification]\label{thm:main}
There exists a countably additive complex-valued measure $\mu:\D\to\mathbb C$
such that $\mu(\cyl(C))=A_p(C)$ for every doubly-principal $C$ if and only if
$p\in[0,1]$. The same existence classification holds if $\mu(\Om)=1$ is
additionally required.
\end{theorem}

A complex-valued measure here takes only finite complex values and is
countably additive on its sigma-algebra. Such a measure has finite total
variation. In particular, for any finite disjoint measurable family,
\begin{equation}\label{eq:variation}
 \sum_i|\mu(E_i)|\le |\mu|(\Om)<\infty.
\end{equation}
This is a standard property of scalar complex measures, not an extra
positivity assumption; compare the bounded-variation framework
in~\cite[Sec.~2.3 and App.~B]{SZ}.

\subsection{Two elementary identities}
For an arbitrary naturally labelled order $G$ of size $k$, let $M(G)$ be
its number of maximal elements and put $w(G)=L(G)+M(G)$. For the empty
order all these counts are zero. The ordinal sum $G\oplus\{u\}\oplus\{v\}$
adds two universal maxima, in that order.

\begin{lemma}[Two caps]\label{lem:caps}
The map $G\mapsto C_G=G\oplus\{u\}\oplus\{v\}$ is a bijection from
naturally labelled $k$-orders to doubly-principal $(k+2)$-orders, and
\begin{equation}\label{eq:caps}
 L(C_G)=L(G)+M(G)+1,\qquad h(C_G)=h(G),\qquad w(G)\ge k.
\end{equation}
\end{lemma}
\begin{proof}
In a finite order a unique maximal element is above every other element.
Natural labelling places it last. Applying this observation twice gives the
unique two-cap decomposition and its converse. The first cap covers exactly
the maximal elements of $G$; the second covers just the first cap. Neither
adds incomparable pairs. Every nonmaximal vertex of $G$ is the lower endpoint
of a cover. Distinct lower endpoints give distinct covers, so
$L(G)\ge k-M(G)$.
\end{proof}

\begin{lemma}[Endpoint product]\label{lem:product}
If $0<r,s<1$ and $r+s=1$, then
\begin{equation}\label{eq:product}
 \sum_{G:\,|G|=k}r^{w(G)}s^{h(G)}
 =F_k(s):=\prod_{j=1}^k(1-s^j)\ge c_s,
 \qquad c_s:=\prod_{j=1}^{\infty}(1-s^j)>0.
\end{equation}
The sum is over distinct naturally labelled orders, without dividing by
automorphisms.
\end{lemma}
\begin{proof}
Independently select each forward edge on $k+1$ ordered vertices with
probability $r$ and take transitive closure. For a prescribed closure $C$,
all its cover edges must occur and edges joining its incomparable pairs
must be absent. Every remaining comparable edge is redundant and its two
choices sum to $r+s=1$. Thus
\begin{equation}\label{eq:closure}
 \Pr(C)=r^{L(C)}s^{h(C)}.
\end{equation}
The event that the last vertex is above all earlier vertices consists of
exactly the orders $G\oplus\{u\}$, and therefore has probability equal
to the sum in~\eqref{eq:product}. Reverse the vertex order and all arrows.
The edge law is unchanged and the event becomes reachability of every later
vertex from the first. Once the first vertex reaches the $j$ vertices
already present, it reaches the next vertex precisely when at least one
of its $j$ incoming edges occurs. These edges have not previously been
exposed, giving the successive factors $1-s^j$.

This is the standard one-sided product used in random-graph post
probabilities~\cite[Sec.~2, Eq.~(2)]{BSW}; the argument is included to fix
the indexing. Finally $\sum_j s^j<\infty$. For sufficiently large $j$,
$\log(1-s^j)\ge-2s^j$, so the tail product has a positive limit. The finite
initial product is positive too.
\end{proof}

\subsection{Variation lower bound and proof of Theorem~\ref{thm:main}}
For $p\notin[0,1]$, put
\begin{equation}\label{eq:normalization}
 a=|p|,\quad b=|1-p|,\quad t=a+b,\quad r=a/t,\quad s=b/t.
\end{equation}
Both $a,b$ are positive. Equality in $|p|+|1-p|\ge1$ holds exactly for
$p\in[0,1]$, so $t>1$. Let
\begin{equation}\label{eq:V}
 V_n(p)=\sum_{\substack{C\text{ doubly principal}\\ |C|=n}}|A_p(C)|.
\end{equation}
By Lemmas~\ref{lem:caps} and~\ref{lem:product},
\begin{align}
 V_{k+2}(p)
 &=a\sum_{|G|=k}a^{w(G)}b^{h(G)}\notag\\
 &=a\sum_{|G|=k}t^{w(G)+h(G)}r^{w(G)}s^{h(G)}\notag\\
 &\ge a t^k F_k(s)\ge a t^k c_s.\label{eq:bound}
\end{align}
Here $w(G)+h(G)\ge k$ and $t>1$; no equality is asserted at the first
inequality. In particular,
\begin{equation}\label{eq:divergence}
 \boxed{V_n(p)\ge |p|\bigl(|p|+|1-p|\bigr)^{n-2}c_s
       \longrightarrow\infty.}
\end{equation}

\begin{proof}[Proof of Theorem~\ref{thm:main}]
For each fixed $n$, the cylinders counted in $V_n$ are disjoint and belong
to $\D$. Any matching complex measure would make $V_n\le|\mu|(\Om)$ for
every $n$ by~\eqref{eq:variation}, contradicting~\eqref{eq:divergence}.
This rules out even unnormalized measures outside $[0,1]$.

For real $0<p<1$, the countable independent-forward-edge probability space
pushes forward under transitive closure to a probability on $(\Om,\F)$.
The closure calculation~\eqref{eq:closure} gives precisely $A_p(C)$ for
every cylinder. Restrict this probability to $\D$. At $p=0$ use the Dirac
probability on the infinite antichain; at $p=1$ use the Dirac probability
on the infinite chain. They have the prescribed endpoint values, with the
zero-power convention stated above, and remain normalized upon restriction.
\end{proof}

\begin{corollary}\label{cor:larger}
The same nonexistence result holds on any sigma-algebra containing $\D$
if the doubly-principal values are retained, in particular on the principal
cylinder sigma-algebra and on $\F$.
\end{corollary}
\begin{proof}
A matching measure on a larger domain would restrict to the measure
excluded by Theorem~\ref{thm:main}.
\end{proof}
The converse inference is invalid: nonexistence on a larger domain alone
would not prove Theorem~\ref{thm:main}.

\section{A separate conditional statement for covariant events}
Let $\mathcal R\subseteq\F$ be the sigma-algebra of measurable events
invariant under order-isomorphism of the infinite histories. A \emph{post}
is an element comparable with every other element. Let $\Pinfty$ be the
event of infinitely many posts, and define
\begin{equation}\label{eq:posts-algebra}
 \Rp=\{E\in\mathcal R:
   \Pinfty^c\subseteq E\text{ or }\Pinfty^c\subseteq E^c\}.
\end{equation}
This is the covariant posts algebra of~\cite[Sec.~V, Eq.~(29)]{DZ}.
It is not the labelled algebra $\D$.

For each finite principal unlabelled shape $H$, let $D_H$ be the event that
the history belongs to $\Pinfty$ and has a post whose inclusive past is
isomorphic to $H$. Let $N(H)$ be the number of distinct naturally labelled
representatives of $H$; labellings related by an automorphism are already
identified in this count. The amplitude $A_p(H)$ is well defined by its
order-invariant exponents. Impose the following additional rule for
\emph{all} principal shapes:
\begin{equation}\tag{T}\label{eq:T}
 \nu(D_H)=\kappa N(H)A_p(H),\qquad
 \kappa\in\mathbb C\text{ independent of }H\text{ and }|H|.
\end{equation}
This is a model assumption about values, not a consequence of covariance.

\begin{theorem}[Conditional common-factor classification]\label{thm:cov}
A countably additive complex-valued measure $\nu$ on $\Rp$ with
$\nu(\Om)=\nu(\Pinfty)=1$ (equivalently, $\Pinfty$ has full
$\nu$-measure) and a finite $\kappa$ satisfying~\eqref{eq:T}
exists if and only if $0<p\le1$.
\end{theorem}
\begin{proof}
First, $D_H$ is measurable: a post with $n-1$ ancestors must have label
$n-1$ in every natural labelling. Its finite past is a finite cylinder
condition and persistence as a post is a countable intersection of finite
conditions. Intersect with the measurable event $\Pinfty$. The result is
isomorphism-invariant and excludes $\Pinfty^c$, so belongs to $\Rp$.

At fixed $n$, different shapes $H$ give disjoint events. Across ranks the
family is laminar: if $|H|<|J|$, then $D_H\cap D_J$ is empty or $D_J$.
Indeed, the finite past $J$ of the later post determines whether an earlier
post of the indicated rank and shape exists inside $J$. Everything beyond
the later post lies above the earlier one as well. The countable family
$\{D_H\}$ covers $\Pinfty$.

If $\kappa=0$, every family member has value zero. Its finite intersections
are empty or family members; inclusion--exclusion makes each finite union
zero. Continuity from below for finite complex measures then gives
$\nu(\Pinfty)=0$, a contradiction. Thus $\kappa\ne0$. This argument uses
laminarity; a union of arbitrary complex-measure zero sets need not have
value zero.

For $p\notin[0,1]$, regroup the labelled sum in~\eqref{eq:V} by unlabelled
doubly-principal shapes. Natural-label multiplicities are exactly $N(H)$,
so the disjoint rank-$n$ events satisfy
\begin{equation}\label{eq:cov-variation}
 \sum_{\substack{[H]\text{ doubly principal}\\ |H|=n}}|\nu(D_H)|
 =|\kappa|\sum_{[H]}N(H)|A_p(H)|=|\kappa|V_n(p).
\end{equation}
This diverges by~\eqref{eq:divergence}, contradicting finite total variation.
At $p=0$, every principal $H$ with $|H|\ge2$ has a cover, so its amplitude
vanishes. Those sizes alone still cover $\Pinfty$; the laminar argument
again contradicts normalization.

For $0<p<1$, classical transitive percolation has infinitely many posts
almost surely, a prior result recorded in~\cite[Sec.~2]{BSW} and attributed
there to Alon et al. Its restriction to $\Rp$ has~\eqref{eq:T} with
\begin{equation}\label{eq:kappa}
 \kappa=c(1-p),\qquad c(q)=\prod_{j\ge1}(1-q^j)>0.
\end{equation}
The finite past has probability $N(H)A_p(H)$. Conditional on that past, the
top remains below every later birth with probability $c(q)$: after it has
reached $j$ vertices, the next birth connects to at least one of them with
probability $1-q^j$. This future factor is independent of the shape of $H$.
The argument is the classical post product~\cite[Theorem~2.1]{BSW}.
At $p=1$ the infinite chain gives the required measure with $\kappa=1$.
\end{proof}

\begin{proposition}[Why all principal shapes are required]\label{prop:dirac}
If~\eqref{eq:T} is required only for doubly-principal $H$, a normalized
measure satisfying it with $\kappa=0$ exists for every numerical $p$.
\end{proposition}
\begin{proof}
Take the infinite ordinal sum
\begin{equation}\label{eq:dirac}
 (A_2\oplus\{z_1\})\oplus(A_2\oplus\{z_2\})\oplus\cdots,
\end{equation}
where $A_2$ is a two-element antichain. Each $z_j$ is a post. The other
elements are incomparable with their partner, hence are not posts.
Deleting $z_j$ from its inclusive past leaves two maximal elements, so no
inclusive post past is doubly principal. The Dirac measure at this history,
restricted to $\Rp$, has $\nu(\Om)=\nu(\Pinfty)=1$ (so $\Pinfty$ has full measure) and assigns zero to every
doubly-principal $D_H$. It satisfies exactly the weakened rule.
\end{proof}

\section{Prior work and the size of the contribution}
Three predecessors locate the mathematical delta.
\emph{Dowker--Zalel}~\cite[Sec.~V, Eqs.~(29)--(31)]{DZ} defines the target
algebra and leaves the doubly-reduced extension test inconclusive. They note that the valency-one sum in their criterion
``in future it could be computed using techniques similar to those used to
obtain the probability of a post.'' The contribution is the explicit bound
\eqref{eq:divergence} and the resulting classification, executing that
suggestion. The individual level sum $V_n$ here is not literally their valency-one sum
$S_{n-r}^{v=1}$.
\emph{Zalel's thesis}~\cite[Ch.~6, Sec.~6.2.3, pp.~89--90]{Zalel}
records the earlier super-reduced formulation and defers the calculation.
The thesis read depth is inherited from the prior source-record lane; this
drafting task did not reopen it, and the load-bearing priority statement is
the published Dowker--Zalel passage above.
\emph{Dowker--Johnston--Surya}~\cite[Sec.~4]{DJS} already establishes
full-labelled complex-percolation nonextension; Surya--Zalel~\cite[Theorem~2.1]{SZ}
states the scalar bounded-variation extension criterion. Failure on those
larger domains does not by itself settle this smaller prescribed-value problem.

The endpoint product is established probability theory, rederived in
Lemma~\ref{lem:product}. No new extension criterion or probability method is
claimed. The conditional statement is additional only under its stated value
rule; arbitrary covariant measures are not classified. Targeted project
priority searches, a citation sweep of the published Dowker--Zalel paper
(INSPIRE and OpenAlex), and independent source checks found no matching proof at
their recorded depths. They are not an exhaustive literature search or a
certificate of global priority. A prior proof of the same prescribed-value
classification would remove the publication delta without invalidating the
argument.

\section{Scope and falsifiers}
\paragraph{Scalar amplitudes versus quantum measures.}
The complex-percolation amplitude is a model choice: no Bell-causality
condition, no operator-valued quantum sequential growth, and no comparison
with the Bell-causal phase results is asserted here. Countable additivity
concerns the complex amplitude measure $\mu$.
The set function $Q(E)=|\mu(E)|^2$ is generally not additive. The theorem
does not apply the probability sum rule to $Q$, nor classify arbitrary
decoherence functionals. Generalized event-convergence prescriptions or
finitely additive assignments are outside its hypotheses.

\paragraph{An infinite-dimensional norm-only escape.}
At $p=(1+i)/2$, let $P_{1/2}$ be classical percolation and define
$m(E)=1_E\in L^2(P_{1/2})$. This standard indicator-valued measure is
countably additive in norm, since the squared norm of a disjoint tail is
its probability. On every finite cylinder,
\[
 \|m(\cyl(C))\|_2^2
 =P_{1/2}(\cyl(C))=|A_p(C)|^2.
\]
Thus matching cylinder norms in a new Hilbert space is possible even where
the prescribed complex measure does not exist. Its vector total variation
is infinite by~\eqref{eq:divergence}. It preserves neither the scalar phases
nor the original off-diagonal decoherence data. It is a scope counterexample,
not an extension of the fixed complex premeasure or a proposed dynamics.

\paragraph{Concrete ways the argument could fail.}
A counterexample to the two-cap bijection or the cover count in
\eqref{eq:caps}, a failure of the endpoint identity~\eqref{eq:product}, or a
matching finite complex measure with unbounded disjoint sums would defeat
the main proof. For Theorem~\ref{thm:cov}, an overlap contradicting laminarity
or a missing natural-label multiplicity would defeat its additional step.
Finite checks below test these encodings; the proofs supply their all-level
justification. Rank-dependent factors and generalized regularizations are
not asserted as theorems in this note.

\paragraph{Physical and philosophical limits.}
The event domain, scalar codomain, countable additivity and matching rule are
mathematical model choices. No observer-access model, empirical prediction or
interpretation of growth as objective passage is supplied. A conclusion about
temporal ontology would require independent physical and philosophical bridge
premises. Neither this theorem nor a finite formal certificate supplies them.

\appendix
\section{Reproducibility and finite formal scope}\label{app:repro}
The accompanying repository contains the exact-arithmetic runner
\path{research/paper-novel/reproduce.py}. From the repository root run
\begin{verbatim}
python3 research/paper-novel/reproduce.py
\end{verbatim}
It records the selected source revision and hashes, runs the accepted
independent finite scripts without changing their outputs, and exits nonzero
on failure. It verifies the two-cap values at $p=(1+i)/2$,
\[
 V_2=\frac1{\sqrt2},\qquad
 V_5=\frac34+\frac{3\sqrt2}{16},\qquad
 V_6=\frac{31}{32}+\frac{55\sqrt2}{64},
\]
along with 407 exact small-cylinder identities through five labels and the
rank-one counterassignment that prevents an incorrect local-renewal extension.
The GLM review suite independently checks 26 finite identities, including
natural-label multiplicities and an older ordinal-block certificate.

No single finite sum exceeding one proves nonextension: an unknown finite
complex measure need not have total variation one. The contradiction is the
\emph{unbounded sequence}~\eqref{eq:divergence}. Enumerated cases do not
establish its infinite-product or all-level steps.

An existing light Lean certificate in \texttt{formal/lean/Certificate.lean}
checks a finite numerical kernel. Its recorded axiom audit does not formalize
sigma-algebras, countable additivity, the positive infinite product, or either
classification theorem. This manuscript does not present those theorems as
machine-verified. AI-assisted derivation and review were used; the mathematical
evidence is the written proof, source correspondence and reproducible checks,
not agreement among agents.

The accompanying \texttt{CLAIM\_LEDGER.md} separates original calculations,
prior results, conditional assumptions and bounded checks, with exact source
locators and read-depth limitations. \texttt{README.md} gives build commands
and the manuscript's review status.

\begin{thebibliography}{9}
\bibitem{DZ} F. Dowker and S. Zalel,
\emph{Observables for cyclic causal set cosmologies},
Class.\ Quantum Grav.\ 40 (2023) 155015, DOI 10.1088/1361-6382/ace149;
arXiv:2212.01149v2, Secs.~III.4 and V.
\href{https://arxiv.org/abs/2212.01149v2}{Versioned preprint}.
\bibitem{Zalel} S. Zalel,
\emph{Discrete random spacetimes: covariance and quantization in growth
dynamics for causal sets}, PhD thesis, Imperial College London (2021),
Ch.~6, Sec.~6.2.3, pp.~89--90.
\href{https://spiral.imperial.ac.uk/entities/publication/dfb9da9a-7d49-4b1a-9f34-7f19b5c9afcf}{Institutional record}.
\bibitem{DJS} F. Dowker, S. Johnston and S. Surya,
\emph{On extending the Quantum Measure}, arXiv:1007.2725v1
(16 July 2010), Sec.~4.
\href{https://arxiv.org/abs/1007.2725v1}{Versioned preprint}.
\bibitem{SZ} S. Surya and S. Zalel,
\emph{A Criterion for Covariance in Complex Sequential Growth Models},
Class.\ Quantum Grav.\ 37 (2020) 195030, DOI 10.1088/1361-6382/ab987f;
arXiv:2003.11311v1, Theorem~2.1 and App.~B.
\href{https://arxiv.org/abs/2003.11311v1}{Versioned preprint}.
\bibitem{BSW} L. Bombelli, I. Seggev and S. Watson,
\emph{A Computation of the Expected Number of Posts in a Finite Random Graph
Order}, arXiv:0809.2258v2 (25 September 2008), Sec.~2, Theorem~2.1,
Eqs.~(1)--(2).
\href{https://arxiv.org/abs/0809.2258v2}{Versioned preprint}.
\end{thebibliography}
\end{document}
